% Results of the checks, 2026-11-27 (Lean v4.33.0, Mathlib v4.33.0).
\paragraph{Build.} \lean{lake build} completes with 1380 jobs, no errors and
no warnings, and prints the axioms of Section~\ref{sec:trust}. A clean build
of the package (with Mathlib taken from its cache) took 8~h~21~min on an
Intel Core i7-8700. \lean{ProofsSweeps16.lean}, the six compiled re-runs on
one thread, takes 8~h~20~min of it. The eight files of
Sections~\ref{sec:statement} and~\ref{sec:def1} build in about 14 seconds.
Later builds reuse the compiled result while the sources it depends on are
unchanged.

\paragraph{Dependency check.} \lean{checks/ProofCone.lean} follows types,
definition bodies and proofs from the two theorems of \lean{Theorem16.lean}
and the seven statements (core Lean and Mathlib are not followed) and lists
every declaration of the package that is never reached. Run in the built
package with \lean{lake env lean checks/ProofCone.lean}, it prints
\begin{code}
declarations of the package: 635; used by the results: 514
not used, written in the source: 0
not used, generated by Lean: 121
\end{code}
The 121 unreached declarations are recursors and other constructions that
Lean generates for every structure and inductive type, derived instances,
generated congruence lemmas, and the two projections \lean{Problem.name} and
\lean{Problem.layer} of fields the results never read. Every declaration
written in the source is used.

\paragraph{Statement cone.} The same traversal restricted to types and
definition bodies gives the declarations whose meaning a reader has to
accept. For the two theorems of \lean{Theorem16.lean} these are the seven
definitions of \lean{Definition1.lean} (61 lines). For Statements~1--4 they
are 41 declarations (257 lines), and for all seven statements 60
declarations (349 lines), three of them proofs that the definitions refer
to; the README lists their line ranges.
